\section{Lean formalization: exact targets and a completion route}
\label{sec:formalization}

\subsection{Two different correctness questions}

There are two tasks. The first is to formalize the mathematical inequalities
using Lean's natural numbers, rationals, and, if convenient for a reference
development, Mathlib's real and complex numbers. The second is to construct
an object-language derivation in \ipra. The first task does not entail the
second. Classical tactics in a reference development must not be silently
exported to the object theory.

The repository provides the syntax \texttt{PR}, \texttt{PTm},
\texttt{Fm}, the intuitionistic proof relation \texttt{Prf}, and the
axioms \texttt{PRAAx}. The required final expression has the shape
\[
 \texttt{Prf PRAAx [] mainFormula},
\]
where \texttt{mainFormula} is the universal closure of
\eqref{eq:integer-matrix} under the threshold antecedent.
Only after constructing that proof may one apply
\texttt{pra\_iff\_native} to obtain its native-calculus counterpart.
The separate bridge file in the archive demonstrates this use with an
\emph{explicitly supplied proof argument}; it does not supply the PNT proof.

\subsection{The implementation sequence}

\begin{longtable}{@{}p{0.19\textwidth}p{0.73\textwidth}@{}}
\toprule
Stage & Concrete proof obligation and endpoint\\
\midrule
\endhead
Arithmetic coding & Define signed integers and positive-denominator rationals;
prove cross-multiplication order, absolute value, finite summation, powers,
factorial, binary length, and bounded search by quantifier-free induction.
Represent every function by a fixed \texttt{PR} expression.\\[3pt]
Primes & Define bounded divisibility and primality, the count $\Pi$,
prime-power detection, and rational approximations to $\psi,\vartheta$.
Verify bounded factorization by the invariant that the product already
extracted times the residual is the original integer. Prove the elementary
unique-factorization identities used in finite Euler products.\\[3pt]
Elementary names & Implement \eqref{eq:log-algorithm}, the finite exponential,
and \eqref{eq:Li-algorithm}. Prove their displayed rational enclosure bounds.
Use explicit positive lower bounds for every reciprocal.\\[3pt]
Finite calculus & Prove interval telescoping with Taylor remainders, the
rectangle identity \eqref{eq:finite-rectangle}, and the explicit square
integral for a simple pole. All mesh sizes and sample accuracies are inputs
or prescribed PR terms.\\[3pt]
Zeta construction & Implement finite Euler summation, with (E3) as its
error budget, and the finite Euler products. Prove (E5)--(E7) with their
cutoffs. This includes the logarithmic derivative identity, not an axiom
asserting it.\\[3pt]
Apartness & Prove the polynomial cosine identity and the finite product
bound $\ge1/2$. Derive the rational-parameter inequalities (A3r), and
the regular expression (A4) on the small-height band.\\[3pt]
Counting formula & Prove the scalar kernel (E8), then its finite-sum
specialization (A5). Apply finite contour cancellation after subtracting the
known pole and obtain (A6) with the stated constants.\\[3pt]
Dyadic estimate & Derive (A7), the two finite-difference inequalities,
(A8), and Theorem~\ref{thm:theta-modulus-established}. Integer roots and
floor divisions are implemented by bounded search.\\[3pt]
Transfer & Instantiate finite Abel summation, use the split at the bounded
integer square root, and prove the factor $17/32$ estimate. Combine with
the error of $\mathcal A$ and clear denominators.\\[3pt]
Export & Universally close the quantifier-free matrix, obtain the actual
\texttt{Prf PRAAx [] mainFormula}, then use the repository's proved
native translation. Audit the dependency tree for added object axioms.\\
\bottomrule
\end{longtable}

This order is a dependency order. It does not require formalizing the
functional equation of $\zeta$, Hadamard factorization, an argument principle
for arbitrary holomorphic functions, an enumeration of zeros, or a general
Tauberian theorem. The main estimate uses one explicitly given pole and
one finite rectangle identity. Existing work on PNT in other proof
assistants or classical Lean blueprints can guide a reference proof, but
their proof-theoretic strength is not automatically that of \ipra.

\subsection{A specific error-budget convention}

It is useful to make the intended compilation mechanically checkable.
An arithmetic lemma carries the precision request and a rational error
allowance as explicit parameters. When a derived expression has a known
Lipschitz multiplier $K$ and uses at most $J$ input samples, evaluate
each sample with error at most
\[
 \frac{2^{-v}}{16(J+1)(K+1)}.
\]
The required number of bits is
$v+4+b(J+1)+b(K+1)$, increased by two if the input interval itself is
rounded. For products, $K$ is obtained from the known magnitude bounds
and the product rule. For the contour lemma, use its more precise
allowance, which accounts for the internal perimeter.

For the last inference of Theorem~\ref{thm:object}, the transfer proof
leaves $15/32$ of the prime-counting allowance, and the Li algorithm
uses less than $19/64$ of the allowance $2^{-v}$. Allocate at most
$2^{-v-4}$ to the auxiliary analytic computations after their displayed
amplification factors. Every cutoff is then a concrete PR expression in
$n,m,v$. The finite invariants establish the remaining integer inequality;
they are not tests of finitely many small instances of that inequality.

\subsection{What the included Lean files establish}

The included source files give precise target definitions, small finite
algebraic proof candidates, and the existing-library bridge. They use no
new axiom called ``PNT'' and do not conceal the analytic input in an
assumption presented as a theorem. They have \emph{not been compiled in
this workspace}: a Lean toolchain was not available. Consequently no
kernel-checking claim is made for these new files. The pinned repository
declares Lean and Mathlib version \texttt{v4.32.0}; the checked build of
that repository was not rerun here.

The completion criterion is stronger than ``the files contain no
\texttt{sorry}''. A finished development must define the PR codes for
the arithmetic programs, identify \texttt{mainFormula} with the actual
matrix, construct its \texttt{Prf PRAAx} derivation, and audit all
dependencies. The present archive supplies a developed mathematical proof
and a concrete route to those objects, as requested, rather than asserting
that this work is already mechanized.

\section{Comparison with the linked strengthenings and inspected drafts}
\label{sec:scope}

The logarithmic integral has two common normalizations:
$\Li_2(x)=\int_2^xdt/\log t$ and the principal-value $\operatorname{li}(x)$.
They differ by $\operatorname{li}(2)$, a fixed computable real. The explicit
arithmetic target in this note uses only $\Li_2$. All asymptotic error
classes asserted here are unchanged by a fixed additive constant;
their numerical thresholds need enlargement if a different normalization
is substituted. We do not transfer a small-argument numerical constant
between the two normalizations without accounting for this difference.

The standard $\log^{-9}T$ argument used here is an established analytic
method; a corresponding classical proof is discussed in the
Kontorovich project blueprint \cite{Kontorovich}. The contribution of this
note is its explicit rational apartness bounds, coarse integer moduli,
finite error accounting, and the specified weak-arithmetic target. The
classical square-root-logarithm estimate requires a stronger zero-free
region; the Vinogradov--Korobov error requires stronger analytic estimates
again. Neither stronger input is proved in this package. The bounds
$O(\sqrt x\log x)$ and Schoenfeld's sharp conditional inequality also
must retain their Riemann-hypothesis assumption. The introductory scope
table is therefore part of the statement of the result, not an omitted
technical qualification. See \cite{Tao,Trudgian} for stronger classical
estimates and their analytic dependencies.

The inspected archive
\texttt{docs/drafts/pnt\_ipra\_constructive\_2026-11-14.zip}
claims a modulus of the shape
\[
 \left|\frac{\psi(x)}x-1\right|\le\frac1q
 \quad\text{for }x\ge2^{\,2^{(2^{20}q^3)^2+1}}.
\]
That stated modulus alone does not yield an error smaller than every
inverse power of $\log x$. The new proof therefore does not assume that
draft as its analytic input. Its Lean artifacts, and the inspected Cauchy
and zeta drafts, explicitly describe themselves as uncompiled outlines.
Their finite-contour ideas are compatible with the independently detailed
rectangle argument here. The evidence file records exactly what was
inspected; it does not claim an exhaustive audit of the whole repository.

There is a related proof-theoretic pitfall in the literature. Avigad
explains that an appeal to conservation of $I\Sigma_1$ over PRA does not
by itself establish the usual quantifier-alternating statement of the
PNT \cite[footnote 14]{Avigad}. Our arithmetic target has a prescribed
modulus and a quantifier-free matrix, and our argument describes its
finite proof directly. We make no blanket conservation claim for arbitrary
analytic theorems.

\section{Reproduction and verification record}

The archive contains the complete LaTeX sources, the compiled PDF,
the finite rational algorithms, Lean outlines, source provenance, and a
machine-readable verification record. The Python checks verify finite
Abel identities exactly, check the cleared-denominator matrix against
the equivalent rational inequality, and compare a small set of logarithm
and Li enclosures with an independent 100-digit calculation. They also
check representative integer prefactor inequalities. These are useful
implementation checks; none is represented as a computational proof of
the universal PNT or of the zero-free strip.

Rebuild the document with two runs of \texttt{pdflatex} from the package
directory. Run \texttt{python3 code/check\_arithmetic.py} for the finite
checks; this requires \texttt{mpmath} only for the independent numerical
comparisons. The arithmetic specification itself uses Python's standard
library. The final PDF was rendered for layout inspection. Detailed
status and commands are in \texttt{README.txt}.
